\documentclass[11pt,a4paper]{article}
\usepackage[utf8]{inputenc}
\usepackage[vietnamese]{babel}
\usepackage{amsmath,amssymb}
\usepackage{geometry}
\geometry{a4paper, total={180mm,260mm}, left=15mm, right=15mm, top=18mm, bottom=18mm}

\usepackage{tikz}
\usetikzlibrary{shapes.geometric, arrows.meta, positioning, calc}
\tikzset{arr/.style={-{Stealth[scale=0.8]}}}
\usepackage{tcolorbox}
\tcbuselibrary{skins,breakable}
\usepackage{fontawesome}
\usepackage{fancyhdr}
\usepackage{enumitem}

\pagestyle{fancy}
\fancyhf{}
\renewcommand{\headrulewidth}{0.5pt}
\renewcommand{\footrulewidth}{0.5pt}
\fancyhead[L]{\small\bfseries\color{violet!85!black}\faBook\ SỔ TAY KIẾN THỨC TOÁN 2}
\fancyhead[R]{\small\bfseries\color{violet!85!black}Lớp Toán Thầy Lĩnh -- loptoanthaylinh.vn}
\fancyfoot[C]{\small\color{gray}-- Sổ tay kiến thức: Bài 04. Hơn, kém nhau bao nhiêu đơn vị --}

\setlength{\parindent}{0pt}
\setlength{\parskip}{0pt}

\begin{document}

% --- BANNER BÀI HỌC ---
\begin{tcolorbox}[
  enhanced,
  colback=violet!8!white,
  colframe=violet!70!black,
  arc=3mm, boxrule=1.2pt,
  left=3mm, right=3mm, top=2mm, bottom=2mm,
  halign=center
]
  {\LARGE\bfseries\color{violet!90!black}\faBookmark\ BÀI 04: HƠN, KÉM NHAU BAO NHIÊU ĐƠN VỊ}\\[1mm]
  {\itshape\footnotesize\color{black!75} Phương pháp tìm phần chênh lệch giữa hai đại lượng và bài toán so sánh trên 1 trang A4}
\end{tcolorbox}

\vspace{-2mm}

\noindent
\begin{minipage}[t]{0.49\textwidth}\vspace{0pt}
\begin{tcolorbox}[
  enhanced,
  colback=violet!4!white,
  colframe=violet!75!black,
  arc=2mm, boxrule=0.8pt,
  left=2mm, right=2mm, top=1.8mm, bottom=1.8mm,
  title={\footnotesize\bfseries\color{white}\faKey\ 1. Quy tắc tìm phần hơn, kém}
]
{\footnotesize
\begin{itemize}\setlength{\itemsep}{1pt}\setlength{\leftskip}{-2.5mm}
  \item Muốn biết số lớn \textbf{hơn} số bé bao nhiêu đơn vị (hoặc số bé \textbf{kém} số lớn bao nhiêu đơn vị), ta luôn thực hiện \textbf{phép tính trừ}:
  $$\mathbf{\text{Phần chênh lệch} = \text{Số lớn} - \text{Số bé}}$$
  \item \textbf{Nhận xét quan trọng:}\\
  Nếu $A$ nhiều hơn $B$ là $5$ đơn vị, thì ngược lại $B$ ít hơn $A$ đúng $5$ đơn vị.\\
  \textit{Ví dụ:} Lan có $12$ bút, Mai có $8$ bút.\\
  Lan hơn Mai: $12 - 8 = 4$ bút.\\
  Mai kém Lan: $12 - 8 = 4$ bút.
\end{itemize}
}
\end{tcolorbox}
\end{minipage}
\hfill
\begin{minipage}[t]{0.49\textwidth}\vspace{0pt}
\begin{tcolorbox}[
  enhanced,
  colback=teal!4!white,
  colframe=teal!75!black,
  arc=2mm, boxrule=0.8pt,
  left=2mm, right=2mm, top=1.8mm, bottom=1.8mm,
  title={\footnotesize\bfseries\color{white}\faSliders\ 2. Kĩ năng giải toán có lời văn}
]
{\footnotesize
\begin{itemize}\setlength{\itemsep}{1pt}\setlength{\leftskip}{-2.5mm}
  \item \textbf{Bước 1: Xác định hai đối tượng:}\\
  Đâu là \textbf{số lớn}, đâu là \textbf{số bé}?
  \item \textbf{Bước 2: Viết câu lời giải:}\\
  Dựa vào câu hỏi: "Số lớn hơn số bé là..." hoặc "Số bé ít hơn số lớn là...".
  \item \textbf{Bước 3: Thiết lập phép tính:}\\
  Luôn lấy: $\text{Số lớn} - \text{Số bé}$.
  \item \textbf{Bước 4: Ghi đáp số:} Kèm đúng đơn vị đo của bài toán.
\end{itemize}
}
\end{tcolorbox}
\end{minipage}

\vspace{-1.5mm}
\begin{tcolorbox}[
  enhanced,
  colback=orange!4!white,
  colframe=orange!80!black,
  arc=2mm, boxrule=0.8pt,
  left=2mm, right=2mm, top=1.5mm, bottom=1.5mm,
  title={\footnotesize\bfseries\color{white}\faSitemap\ 3. Sơ đồ đoạn thẳng so sánh phần hơn và phần kém}
]
\centering
\begin{tikzpicture}[x=1cm, y=0.8cm]
  % Đoạn thẳng số lớn
  \node[left, font=\scriptsize\bfseries] at (0, 1) {Số lớn:};
  \draw[thick, violet!80!black, fill=violet!15] (0, 0.8) rectangle (5, 1.2);
  \node[font=\scriptsize\bfseries] at (2.5, 1) {Số lớn ($15$)};

  % Đoạn thẳng số bé
  \node[left, font=\scriptsize\bfseries] at (0, 0) {Số bé:};
  \draw[thick, teal!80!black, fill=teal!15] (0, -0.2) rectangle (3.2, 0.2);
  \node[font=\scriptsize\bfseries] at (1.6, 0) {Số bé ($9$)};

  % Phần chênh lệch nét đứt
  \draw[dashed, red!80!black, line width=0.8pt] (3.2, 0.2) -- (3.2, 1.2);
  \draw[dashed, red!80!black, line width=0.8pt] (5, 0.8) -- (5, -0.2);
  \draw[<->, thick, red!85!black] (3.2, 0) -- (5, 0) node[midway, below=2pt, font=\scriptsize\bfseries] {Hơn/Kém: $15 - 9 = 6$};
\end{tikzpicture}
\end{tcolorbox}

\vspace{-1.5mm}
\noindent
\begin{minipage}[t]{0.49\textwidth}\vspace{0pt}
\begin{tcolorbox}[
  enhanced,
  colback=green!4!white,
  colframe=green!60!black,
  arc=2mm, boxrule=0.8pt,
  left=2mm, right=2mm, top=1.8mm, bottom=1.8mm,
  title={\footnotesize\bfseries\color{white}\faCheckCircle\ 4. Ví dụ thực hành chuẩn mực}
]
{\footnotesize
\textbf{Bài toán mẫu:}\\
Bạn An câu được $16$ con cá, bạn Bình câu được $9$ con cá. Hỏi bạn An câu được nhiều hơn bạn Bình bao nhiêu con cá?\\
\textit{Bài giải:}\\
Bạn An câu được nhiều hơn bạn Bình số con cá là:
$$16 - 9 = 7\text{ (con cá)}$$
Đáp số: $7$ con cá.
}
\end{tcolorbox}
\end{minipage}
\hfill
\begin{minipage}[t]{0.49\textwidth}\vspace{0pt}
\begin{tcolorbox}[
  enhanced,
  colback=red!4!white,
  colframe=red!75!black,
  arc=2mm, boxrule=0.8pt,
  left=2mm, right=2mm, top=1.8mm, bottom=1.8mm,
  title={\footnotesize\bfseries\color{white}\faExclamationTriangle\ 5. Mẹo tránh bẫy trắc nghiệm}
]
{\footnotesize
\begin{itemize}\setlength{\itemsep}{1pt}\setlength{\leftskip}{-2.5mm}
  \item \textbf{Bẫy thấy chữ "nhiều hơn" vội làm tính cộng:}\\
  Khi bài toán hỏi \textbf{"hơn bao nhiêu"} hoặc \textbf{"kém bao nhiêu"}, ta \textbf{bắt buộc làm phép trừ} ($\text{Lớn} - \text{Bé}$).
  \item \textbf{Quy tắc vàng:}\\
  Hỏi "nhiều hơn bao nhiêu": Lấy trừ.\\
  Hỏi "ít hơn bao nhiêu": Cũng lấy trừ!
  \item Luôn lấy số lớn trừ số bé, không lấy số bé trừ số lớn.
\end{itemize}
}
\end{tcolorbox}
\end{minipage}


\end{document}
